\documentclass{article}
\usepackage[T1]{fontenc}
\usepackage{lmodern}
\usepackage{graphicx}
\usepackage{amsmath,amssymb}
\usepackage[a4paper,margin=2cm]{geometry}
\usepackage[absolute,overlay]{textpos}
\setlength{\TPHorizModule}{1mm}
\setlength{\TPVertModule}{1mm}
\usepackage[indonesian]{babel}
\usepackage{hyperref}
\setlength{\emergencystretch}{2em}
% unit: Halaman 21

\begin{document}

% === Halaman 21 (llm) ===

\begin{itemize}
    \item Secara umum dapat disimpulkan bahwa RSA hanya aman jika $n$ cukup besar.
    \item Jika panjang $n$ hanya 256 bit atau kurang, ia dapat difaktorkan dalam beberapa jam saja dengan sebuah komputer $PC$ dan program yang tersedia secara bebas.
    \item Jika panjang $n$ adalah 512 bit atau kurang, ia dapat difaktorkan dengan beberapa ratus komputer.
    \item Saat ini panjang kunci RSA yang aman adalah jika $n$ di atas 1024 bit.
\end{itemize}

\end{document}
